\begin{lemma}[A reversed path extends to a normal maximum tree]
Let $(T,e_0,\varphi)$ be a globally maximum $k$-triple, assume that $T$
is fixed-pair maximal, and write its first two vertices as $y_0,y_1$.
Suppose that $\alpha\in\overline\varphi(y_0)$ and
$\beta\in\overline\varphi(y_1)$, and let $P$ be a simple
$(\alpha,\beta)$-incidence path from $y_0$ to $y_1$.  From any reversed
Kempe-path certificate $C$ for $P$, one obtains a globally maximum Tashkinov
tree $T'$ and a normal witness on $T'$ whose height is
$\operatorname{length}(P)/2+1$.
\end{lemma}
\begin{proof}
Let $U$ and $q$ be the tree and closing edge supplied by
\noderef{ScheideReversedPathCertificate}.  Its fields say that $U$ is a
Tashkinov tree and has length
$m=\operatorname{length}(P)/2\geq2$.  Apply
\noderef{ScheideMaximalExtendsEveryTree} to the $k$-triple contained in
\noderef{GlobalMaximumTriple}, the assumed fixed-pair maximal tree $T$,
and $U$.  We obtain a Tashkinov tree $T'$ such that
\[
 V(T')=V(T),\qquad U\text{ is an ordered initial segment of }T'. \tag{1}
\]
By \noderef{TreeSequenceInitialSegment}, (1) gives
$m\leq T'.length$ and exact equality of every vertex entry through index
$m$ and every edge entry below $m$.  The cardinality equality in
\noderef{GlobalMaximumTriple}, together with $V(T')=V(T)$, shows that
$(T',e_0,\varphi)$ is again globally maximum.

We construct the normal witness field by field.  Set its height to
$m+1$ and its last index to the index $m$ of $T'$.  The certificate gives
$m\geq2$, while (1) gives $m\leq T'.length$, so the required height bounds
hold.  We next verify the two missing-color fields explicitly.  The
certificate equalities are
\[
 U.vertex(0)=T.vertex(0)=y_0,
 \qquad U.vertex(1)=T.vertex(1)=y_1.
\]
The vertex-entry clause of (1), first at index $0$ and then at index $1$,
gives
\[
 T'.vertex(0)=U.vertex(0),
 \qquad T'.vertex(1)=U.vertex(1).
\]
Combining these equalities yields $T'.vertex(0)=y_0$ and
$T'.vertex(1)=y_1$.  Substitution in the two explicit hypotheses
$\alpha\in\overline\varphi(y_0)$ and
$\beta\in\overline\varphi(y_1)$ therefore gives, respectively,
$\alpha\in\overline\varphi(T'.vertex(0))$ and
$\beta\in\overline\varphi(T'.vertex(1))$.  These are exactly the
`alpha_missing' and `beta_missing' fields of the normal witness.  Choose
$q$ as the closing edge.  The certificate's closing-end equation,
transported entrywise through (1), is
\[
 G.\operatorname{ends}(q)
   =\{T'.vertex(m),T'.vertex(0)\}.                              \tag{2}
\]

Use the reversed walk $P^{\rm rev}$ as the closing walk.  It starts at
the second vertex and ends at the first, because $P$ has the opposite
endpoints and the first two entries of $U,T'$ agree.  Reversal preserves
simplicity, support, and length.  Since $m=\operatorname{length}(P)/2$
and an incidence path between vertex-nodes has even length, this length is
exactly $2m=2((m+1)-1)$.

The certificate's vertex-support equivalence says that the vertex-nodes
of $P$, hence also of $P^{\rm rev}$, are exactly the $m+1$ vertex entries
of $U$.  Entrywise equality in (1) identifies these with precisely the
first $m+1$ entries of $T'$, giving `closingPath_vertices'.  Likewise,
the certificate's edge-support equivalence says that its edge-nodes are
exactly $q$ together with the entries $U.edge(j)$ for
$1\leq j\leq m-1$.  Under (1), these are exactly $q$ together with the
entries $T'.edge(j)$ satisfying
$1\leq j$ and $j+2\leq m+1$; this proves both directions of
`closingPath_edges', including its index bounds.

Finally, for every such $j$, the certificate gives
\[
 G.\operatorname{ends}(U.edge(j))
   =\{U.vertex(j),U.vertex(j+1)\}.
\]
Transporting the three exact entries through (1) gives the corresponding
`trunk_edge_ends' equation for $T'$.  Thus (2), the reversed-walk facts,
the two exact support equivalences, and the trunk equations establish
every field of \noderef{NormalGlobalMaximumTriple}.  Its height is
$m+1=\operatorname{length}(P)/2+1$, as claimed.
\end{proof}
